\paragraph{ZG 的素数 J-分式、黄金 Cantor 地址与 Archimedean/\texorpdfstring{$2^L$}{2^L}-进规模分离}
在 \(ZG\) 的矩阵 Euler 递推、自然密度的前缀单调极限以及 H.161 的 \(2^L\)-进地址像同胚都已固定之后，还可把算术密度与非阿基米德地址几何进一步压缩为下列三条结果：前缀密度具有显式的 prime \(J\)-分式与无穷乘积表示；相应地址像恰是黄金均值子移位的 Cantor 型实现，并具有 \(\log\varphi/\log B\) 的精确维数；从而 \(ZG\) 在 Archimedean 计数侧的正密度与其在 \(2^L\)-进地址侧的严格亚一维几何之间形成一条明确的尺度分离。

\begin{theorem}[ZG 自然密度的素数 \texorpdfstring{$J$}{J}-分式、无穷乘积与定量正性]\label{thm:xi-time-part62dgd-zg-density-prime-jfraction-positive}
\leanverified{paper\_xi\_time\_part62dgd\_zg\_density\_prime\_jfraction\_positive}
记
\[
\mathbf 1:=(1,1)^{\mathsf T},
\qquad
e_1:=(1,0)^{\mathsf T},
\]
并定义
\[
\binom{A_n}{B_n}
:=
\left(
\prod_{i=1}^{n}
\begin{pmatrix}
\dfrac{p_i}{p_i+1} & \dfrac{p_i}{p_i+1}\\[1ex]
\dfrac{1}{p_i+1} & 0
\end{pmatrix}
\right)e_1,
\qquad
t_n:=A_n+B_n,
\qquad
r_n:=\frac{B_n}{A_n},
\]
且约定 \(r_0:=0\)。则 \(\delta_{\mathrm{ZG}}=\mathrm{dens}(ZG)\) 满足：
\begin{enumerate}
  \item 对每个 \(n\ge 1\)，有递推
  \[
  r_n=\frac{1}{p_n(1+r_{n-1})},
  \]
  从而得到 \(J\)-分式展开
  \[
  r_n=
  \cfrac{1}{p_n\!\left(1+\cfrac{1}{p_{n-1}\!\left(1+\cfrac{1}{\ddots\left(1+\frac1{p_1}\right)}\right)}\right)}.
  \]
  \item 对每个 \(n\ge 1\)，有
  \[
  t_n
  =
  t_{n-1}\left(1-\frac{r_{n-1}}{(p_n+1)(1+r_{n-1})}\right).
  \]
  因而
  \[
  \delta_{\mathrm{ZG}}
  =
  \frac1{\zeta(2)}\lim_{n\to\infty}t_n
  =
  \frac1{\zeta(2)}
  \prod_{n\ge 2}
  \left(1-\frac{r_{n-1}}{(p_n+1)(1+r_{n-1})}\right).
  \]
  \item 若记
  \[
  \delta_n:=\frac{t_n}{\zeta(2)},
  \]
  则 \(\delta_n\downarrow \delta_{\mathrm{ZG}}\in(0,1)\)，并且对每个 \(n\ge 1\) 有显式尾估计
  \[
  0<\delta_n-\delta_{\mathrm{ZG}}
  \le
  \delta_n\sum_{j\ge n+1}\frac{1}{p_{j-1}(p_j+1)}.
  \]
\end{enumerate}
\end{theorem}
\begin{proof}
命题 \ref{prop:xi-zg-prefix-recursion-monotone-tail-certificate} 中的归一化前缀量 \((a_n,b_n)\) 满足
\[
(a_0,b_0)=(1,0),
\]
\[
a_n=(a_{n-1}+b_{n-1})\frac{p_n}{p_n+1},
\qquad
b_n=\frac{a_{n-1}}{p_n+1}.
\]
而这里按矩阵定义的 \((A_n,B_n)\) 由乘法递推也满足同一组初值与同一递推式。故
\[
A_n=a_n,
\qquad
B_n=b_n,
\qquad
t_n=A_n+B_n=u_n.
\]
于是
\[
r_n=\frac{B_n}{A_n}
=
\frac{1}{p_n}\frac{A_{n-1}}{A_{n-1}+B_{n-1}}
=
\frac{1}{p_n(1+r_{n-1})},
\]
第一条得证；递归展开即给出所述 \(J\)-分式。

再由
\[
t_n=A_n+B_n
=
\frac{p_n}{p_n+1}(A_{n-1}+B_{n-1})
+\frac{1}{p_n+1}A_{n-1},
\]
以及
\[
A_{n-1}=\frac{t_{n-1}}{1+r_{n-1}},
\]
可得
\[
t_n
=
t_{n-1}\frac{p_n+\frac{1}{1+r_{n-1}}}{p_n+1}
=
t_{n-1}\left(1-\frac{r_{n-1}}{(p_n+1)(1+r_{n-1})}\right).
\]
这就得到第二条的递推式。又由 \(t_1=1\) 与命题
\ref{prop:xi-zg-prefix-recursion-monotone-tail-certificate}
中的
\[
\delta_{\mathrm{ZG}}=\frac{1}{\zeta(2)}\lim_{n\to\infty}t_n
\]
可知对 \(n\ge 2\) 迭代上式便得到所示无穷乘积表示。

现证第三条。由 \(r_0=0\) 与第一条递推归纳得到
\[
0<r_n\le \frac1{p_n}\qquad (n\ge 1).
\]
记
\[
c_j:=\frac{r_{j-1}}{(p_j+1)(1+r_{j-1})}\qquad (j\ge 2).
\]
则
\[
0<c_j\le \frac{1}{p_{j-1}(p_j+1)}.
\]
由第二条可得
\[
t_\infty
=
t_n\prod_{j\ge n+1}(1-c_j),
\]
从而
\[
0<t_n-t_\infty
=
t_n\left(1-\prod_{j\ge n+1}(1-c_j)\right).
\]
对正数列 \((c_j)\) 应用不等式
\[
1-\prod_{j\ge n+1}(1-c_j)\le \sum_{j\ge n+1}c_j
\]
便得
\[
0<t_n-t_\infty
\le
t_n\sum_{j\ge n+1}\frac{1}{p_{j-1}(p_j+1)}.
\]
两边同时乘以 \(\zeta(2)^{-1}\)，即得
\[
0<\delta_n-\delta_{\mathrm{ZG}}
\le
\delta_n\sum_{j\ge n+1}\frac{1}{p_{j-1}(p_j+1)}.
\]
又命题 \ref{prop:xi-zg-prefix-recursion-monotone-tail-certificate} 已证明 \(t_n\) 单调下降并收敛到 \(\zeta(2)\delta_{\mathrm{ZG}}\)，故 \(\delta_n\downarrow \delta_{\mathrm{ZG}}\in(0,1)\)。证毕。
\end{proof}

\begin{theorem}[ZG 的 \texorpdfstring{$2^L$}{2^L}-进地址像是黄金 Cantor 集，且维数为 \texorpdfstring{$\log\varphi/\log B$}{log phi / log B}]\label{thm:xi-time-part62dgd-zg-address-golden-cantor-dimension}
设
\[
\Sigma_\varphi
:=
\{\,z=(z_k)_{k\ge 1}\in \{0,1\}^{\NN}: z_kz_{k+1}=0\ \forall k\,\},
\]
固定 \(T:=T_2(E_{\min})\) 中一个非零向量 \(v\)，并令 \(B:=2^L\ge 2\)。定义地址映射
\[
X_B:\Sigma_\varphi\longrightarrow T,
\qquad
X_B(z):=\sum_{k\ge 1}z_kB^{k-1}v,
\]
以及像集
\[
\mathcal X_B:=X_B(\Sigma_\varphi).
\]
则成立：
\begin{enumerate}
  \item \(X_B\) 是 \(\Sigma_\varphi\) 到 \(\mathcal X_B\) 的拓扑同胚。若定义
  \[
  \mathcal T:=X_B\circ \sigma\circ X_B^{-1},
  \]
  其中 \(\sigma\) 为黄金均值子移位上的左移，则 \((\mathcal X_B,\mathcal T)\) 与 \((\Sigma_\varphi,\sigma)\) 拓扑共轭。
  \item 在由最长公共前缀等价地由模 \(B^mT\) 同余诱导的 \(B\)-进超度量下，
  \[
  \dim_{\mathrm H}(\mathcal X_B)
  =
  \underline{\dim}_{\mathrm M}(\mathcal X_B)
  =
  \overline{\dim}_{\mathrm M}(\mathcal X_B)
  =
  \frac{\log\varphi}{\log B}.
  \]
  \item 若记 \(N_m(\mathcal X_B)\) 为 \(\mathcal X_B\) 在尺度 \(B^{-m}\) 上的最少球覆盖数，则
  \[
  N_m(\mathcal X_B)=F_{m+2},
  \]
  并且对
  \[
  s:=\frac{\log\varphi}{\log B}
  \]
  有
  \[
  \lim_{m\to\infty}N_m(\mathcal X_B)\,B^{-ms}
  =
  \frac{\varphi^2}{\sqrt5}.
  \]
\end{enumerate}
\end{theorem}
\begin{proof}
取 \(E=\{0,1\}\) 且 \(\mathrm{code}(0)=0\)、\(\mathrm{code}(1)=1\)。由于 \(B=2^L>1\)，定理
\ref{thm:xi-time-part62d-hologram-image-cantor-homeomorphism}
适用于该二元字母表，并给出地址映射
\[
X:E^{\NN}\to \Lambda_{E,v}
\]
是到其像的拓扑同胚。由于 \(\Sigma_\varphi\subset E^{\NN}\) 是闭的移位不变子集，故其限制
\[
X_B:=X|_{\Sigma_\varphi}
:
\Sigma_\varphi\to \mathcal X_B
\]
仍是拓扑同胚。由此定义
\[
\mathcal T:=X_B\circ \sigma\circ X_B^{-1}
\]
便立即得到第一条。

对第二、三条，记 \([\mathbf a]_\varphi\subset \Sigma_\varphi\) 为长度 \(m\) 的合法前缀 \(\mathbf a=(a_1,\dots,a_m)\) 所对应的圆柱集。由定理
\ref{thm:xi-time-part62d-hologram-image-cantor-homeomorphism}
的前缀同余性质，所有 \(z\in[\mathbf a]_\varphi\) 的像都满足
\[
X_B(z)\equiv \sum_{k=1}^{m}a_kB^{k-1}v\pmod{B^mT}.
\]
不同合法前缀对应不同的模 \(B^mT\) 余类，因此在 \(\mathcal X_B\) 上的长度 \(m\) 圆柱集恰给出一族两两不交的半径 \(B^{-m}\) 的闭开球。于是尺度 \(B^{-m}\) 的最少覆盖数恰等于长度 \(m\) 的合法前缀个数。该个数即为标准 Fibonacci 计数
\[
N_m(\mathcal X_B)=F_{m+2}.
\]

由此立得
\[
\underline{\dim}_{\mathrm M}(\mathcal X_B)
=
\overline{\dim}_{\mathrm M}(\mathcal X_B)
=
\lim_{m\to\infty}\frac{\log F_{m+2}}{m\log B}
=
\frac{\log\varphi}{\log B}.
\]
对 Hausdorff 维数，若 \(s>\log\varphi/\log B\)，则长度 \(m\) 圆柱集构成一个直径 \(B^{-m}\) 的覆盖，并满足
\[
F_{m+2}B^{-ms}\longrightarrow 0,
\]
故 \(\mathcal H^s(\mathcal X_B)=0\)。反之，若 \(s<\log\varphi/\log B\)，则任何直径不超过 \(B^{-m}\) 的覆盖都至少需要 \(F_{m+2}\) 个集合，因而其 \(s\)-维代价至少为
\[
F_{m+2}B^{-ms}\longrightarrow \infty.
\]
于是 \(\mathcal H^s(\mathcal X_B)=\infty\)。因此
\[
\dim_{\mathrm H}(\mathcal X_B)=\frac{\log\varphi}{\log B},
\]
第二条得证。

最后，由 Binet 公式
\[
F_{m+2}=\frac{\varphi^{m+2}-(-\varphi^{-1})^{m+2}}{\sqrt5}
\]
可得
\[
F_{m+2}\varphi^{-m}\longrightarrow \frac{\varphi^2}{\sqrt5}.
\]
而
\[
B^{-ms}=\varphi^{-m},
\]
故
\[
N_m(\mathcal X_B)\,B^{-ms}
=
F_{m+2}\varphi^{-m}
\longrightarrow
\frac{\varphi^2}{\sqrt5}.
\]
第三条成立。证毕。
\end{proof}

\begin{corollary}[ZG 的 Archimedean 正密度与 \texorpdfstring{$2^L$}{2^L}-进亚一维地址几何严格分离]\label{cor:xi-time-part62dgd-zg-archimedean-nonarchimedean-scale-separation}
沿用定理 \ref{thm:xi-time-part62dgd-zg-density-prime-jfraction-positive} 与定理 \ref{thm:xi-time-part62dgd-zg-address-golden-cantor-dimension} 的记号。则 \(ZG\) 同时满足
\[
0<\delta_{\mathrm{ZG}}<1,
\]
并且其 \(2^L\)-进地址像 \(\mathcal X_B\) 满足
\[
\dim_{\mathrm H}(\mathcal X_B)
=
\frac{\log\varphi}{\log B}
<1.
\]
因此，\(ZG\) 在整数轴上的自然计数中是正密度对象，而在 \(2^L\)-进地址侧却形成一个严格亚一维的 Cantor 集。
\end{corollary}
\begin{proof}
第一式由定理 \ref{thm:xi-zg-abel-residue-log-density} 与定理
\ref{thm:xi-time-part62dgd-zg-density-prime-jfraction-positive}
给出。第二式是定理
\ref{thm:xi-time-part62dgd-zg-address-golden-cantor-dimension}
的直接结论。又因 \(B=2^L\ge 2>\varphi\)，故
\[
\frac{\log\varphi}{\log B}<1.
\]
证毕。
\end{proof}

\endinput
